% Theorem-Licensed Typestate — A Verification Method for Hardware State Machines
% Target: FMCAD 2026 or MEMOCODE 2026
% Copyright 2026 John Curley. Licensed under CC0 1.0.

\documentclass[conference]{IEEEtran}
\usepackage{cite}
\usepackage{amsmath,amssymb,amsfonts}
\usepackage{algorithmic}
\usepackage{graphicx}
\usepackage{textcomp}
\usepackage{xcolor}
\usepackage{tikz}
\usepackage{booktabs}
\usepackage{multirow}

\usetikzlibrary{positioning,shapes,arrows,chains,fit,backgrounds}

\definecolor{cellblue}{RGB}{220,230,241}
\definecolor{cellred}{RGB}{252,228,214}

\begin{document}

\title{Theorem-Licensed Typestate:\\A Verification Method for Hardware State Machines}

\author{
\IEEEauthorblockN{John Curley}
\IEEEauthorblockA{Sovereign Processor Unit\\
Auckland, New Zealand\\
john@deft.work}
}

\maketitle

\begin{abstract}
We present a three-layer verification discipline for hardware state machines
that synthesises typestate analysis (Strom \& Yemini, 1986), design-by-contract
with algebraic theorems, and structurally independent oracle testing.  The
method encodes algebraic theorems as register bits, guards transitions by the
theorem's hypotheses, and refuses exactly where the hypothesis fails---error is
a terminal state, never silent drift.  Three independent verification layers
backstop the design: a structurally different exact oracle, bit-exact trace
equivalence, and poison-proof fault injection where every guard is proven to
leave the destination register untouched on fault.  A real bug ledger from FPGA
development validates the method: each of three layers independently caught a
bug the others would have missed, including a 50.5\%-failure-rate
cross-catalog truncation caught by typestate analysis before any RTL existed.
We provide five case studies spanning arithmetic accelerators and a protocol
state machine, a measured head-to-head comparison with SystemVerilog
Assertions, and synthesis metrics for all subsystems on a Gowin GW5A-25 FPGA.
\end{abstract}

\begin{IEEEkeywords}
typestate, hardware verification, design-by-contract, oracle testing,
fault injection, FPGA, formal methods
\end{IEEEkeywords}

%=============================================================================
\section{Introduction}
%=============================================================================

A hardware state machine where the states are algebraic theorems encoded as
register bits, transitions are guarded by the theorem's hypotheses, and the
machine refuses exactly where the theorem's hypothesis fails.

It is not a new paradigm.  It is \textbf{typestate}~\cite{strom86} applied to
RTL, \textbf{design-by-contract}~\cite{meyer92} where the contracts are
algebraic theorems rather than arbitrary predicates, and the \textbf{oracle
discipline} inherited from cryptography (where the reference implementation
must not share the implementation's arithmetic model, or it verifies nothing).
The contribution is the synthesis---nobody else packages these three together
for hardware verification, and the bug ledger (Section~\ref{sec:bugs}) shows
why each layer independently earns its keep.

Three verification layers backstop the design:
\begin{enumerate}
\item An independent exact oracle using structurally different arithmetic
  (Layer~1).
\item Bit-exact trace equivalence between oracle and implementation (Layer~2).
\item Poison-proof fault injection where every guard is proven to leave the
  destination register bit-identically untouched on fault (Layer~3).
\end{enumerate}

%=============================================================================
\section{The Bug Ledger}\label{sec:bugs}
%=============================================================================

Every bug in this ledger is attributable to one specific layer; no layer is
redundant.  The three bugs were found in real FPGA development of the SPU-13
rational-field processor~\cite{spu13}, not constructed as examples.

\begin{table}[h]
\centering
\caption{Three bugs, three independent layers}
\label{tab:bugs}
\begin{tabular}{p{2.2cm}p{1.4cm}p{2.8cm}}
\toprule
\textbf{Bug} & \textbf{Layer} & \textbf{How caught} \\
\midrule
CATMIX --- 50.5\% mixed-catalog chain corruption
& Typestate
& Doubling-theorem assert caught before any RTL existed.  Fixed by 4-state
  $\varphi$-plane typestate. \\
\midrule
Thirds $/3$ silent truncation
& Oracle
& Python \texttt{Fraction} vs.\ RTL Q12 fixed-point; structurally different
  arithmetic revealed the silent floor-division. \\
\midrule
Batch-inverter stale read
& Golden-vector
& Ordering-adversarial unit-LAST vectors in a committed \texttt{.mem} file
  exposed a shared-multiplier mux isolation bug. \\
\bottomrule
\end{tabular}
\end{table}

The CATMIX argument is the method's thesis in one sentence: \textbf{the bug
was common (50.5\% hit rate on mixed-catalog chains) but the test was
unwritable without the theory forcing the hypothesis check.}  The spec asserted
mixing was safe, so no test plan would ever have written a mixed-catalog test.
Typestate analysis generates the tests that the spec's own claims would have
suppressed.

\subsection{Why the Oracle Must Be Structurally Different}

If the oracle shares the implementation's arithmetic model (same fixed-width
integers, same rounding discipline), agreement proves nothing---both could
share the same silent error.  The SPU oracles use Python \texttt{Fraction}
(arbitrary-precision rationals) or exact integer arithmetic with no fixed-width
constraints.  The RTL uses 16-bit Q12 fixed-point with truncating division.
When they match bit-for-bit, it is because the arithmetic genuinely closes---not
because they share a flawed model.

This is standard practice in cryptography (test vectors from an independent
implementation) but rare in hardware verification, where gold models typically
reuse the same word width and rounding.

%=============================================================================
\section{The Three-Layer Verification Discipline}
%=============================================================================

\subsection{Layer 1 --- Independent Exact Oracle}

An oracle that computes the same function using a structurally different
arithmetic model.  The oracle is the ground truth; the implementation is
correct iff it matches the oracle bit-for-bit on all tested vectors.

\begin{table}[h]
\centering
\caption{Oracle-implementation structural differences}
\label{tab:oracles}
\begin{tabular}{p{1.8cm}p{2.2cm}p{2.2cm}}
\toprule
\textbf{Subsystem} & \textbf{Oracle} & \textbf{Implementation} \\
\midrule
ROTC thirds & Python arbitrary-precision ints & RTL 4-bit tagged exponents +
Q12 fixed-point \\
IROTC & Python \texttt{Fraction} & VM integer $\mathbb{Z}[\varphi]$ with
\texttt{>{}>{}>1} truncation \\
Lucas MAC & Python \texttt{int} mod $L_p$ & RTL Barrett reduction \\
Batch inv. & Python tower model & RTL shared-multiplier mux \\
\bottomrule
\end{tabular}
\end{table}

\subsection{Layer 2 --- Bit-Exact Trace Equivalence}

For every oracle-verified operation, the VM and RTL must produce identical
register state at instruction granularity---per-operation end states, plus
pinned instruction latency where the timing is architectural (the IROTC
engine's fixed 13-cycle slot is a testbench assertion).  This catches
implementation divergence
(permutation order, writeback timing, sign-extension width) that the oracle's
functional correctness test would miss.

The ROTC trace equivalence suite exercises all 36 angles through both rotor
datapaths---336 bit-exact checks.  The IROTC trace suite exercises all 60
indices $\times$ both catalogs on tagged inputs---9 explicit checks driving 360
component comparisons.

\subsection{Layer 3 --- Poison-Proof Fault Injection}

Every guard that can reject an operation must be proven to leave the
destination register bit-identically untouched.  A fault that fires but
corrupts state is worse than no fault at all---it creates silent data
corruption behind a ``safe'' error code.

\begin{table}[h]
\centering
\caption{Poison-proof fault injection proofs}
\label{tab:poison}
\begin{tabular}{p{2.2cm}p{2.4cm}p{1.6cm}}
\toprule
\textbf{Guard} & \textbf{Test file} & \textbf{Checks} \\
\midrule
ROTC bad-angle gate & \texttt{test\_rotc\_bad\_angle.py} & poison holds \\
IROTC BADIDX/UNTAGGED/CATMIX & \texttt{test\_irotc\_poison.py} & 14 \\
SPI UNKNOWN\_CMD/CRC & \texttt{spi\_protocol\_oracle.py} & 8 \\
\bottomrule
\end{tabular}
\end{table}

All faults are \textbf{terminal} in the typestate lattice: FAULT, OVERFLOW,
DIVERGED, and SINGULAR states have no outgoing transitions.  Error is not
drift---it is a terminal state.  Recovery is always explicit (SCALE2
re-conditioning, RESET, DRAIN), never automatic.

%=============================================================================
\section{The Typestate Lattice}\label{sec:lattice}
%=============================================================================

The $\varphi$-plane typestate is a proper join-semilattice, not a set of
independent boolean flags:

\begin{figure}[h]
\centering
\includegraphics[width=\columnwidth]{fig_synthesis_comparison.pdf}
\caption{Probe-top resource utilization on Tang Primer 25K (GW5A-25A),
parsed from the named nextpnr logs at figure-generation time
(\texttt{tools/gen\_paper\_figures.py}---no hand-embedded numbers).
LUT4 and DFF per probe, Fmax annotated; solid fill = silicon-verified.
Inset: typestate vs.\ SVA-style guard, identical post-synth cell counts.}
\label{fig:synthesis}
\end{figure}

\begin{table}[h]
\centering
\caption{Synthesis comparison: typestate vs.\ SVA-style guards}
\label{tab:sva}
\begin{tabular}{lcc}
\toprule
\textbf{Metric} & \textbf{Typestate} & \textbf{SVA-style} \\
\midrule
RTL lines & 121 & 150 (+24\%) \\
Post-synth cells & 52 & 52 (identical) \\
DFF (state bits) & 5 & 5 \\
LUT4 & 7 & 8 \\
Lattice join operator & 15 lines, reusable & duplicated per opcode \\
Cross-plane hazard & structural (CATMIX) & manual enumeration \\
Guards per new opcode & 0 (reuses join) & re-enumerate per opcode \\
\bottomrule
\end{tabular}
\end{table}

With correct specification semantics the two versions synthesise to
\emph{identical} cell counts---the lattice costs nothing in area.  The
advantage is therefore purely structural: the join operator automatically
catches CATMIX-class cross-plane hazards that the boolean-flag version
requires the designer to enumerate by hand, and a new opcode needs zero new
guards in the lattice version versus a fresh enumeration in the SVA-style
version.

%=============================================================================
\section{Related Work}
%=============================================================================

Typestate was introduced by Strom and Yemini~\cite{strom86} as a programming
language concept for tracking the allowable operation sequences on an object.
Recent work has applied typestate to embedded Rust~\cite{coconut24} and to
session types, but these are compile-time disciplines operating on software
types, not runtime register-level hardware state.

SystemVerilog Assertions (SVA)~\cite{sva1800} is the dominant hardware
assertion language.  SVA assertions are boolean predicates over signals; they
lack the lattice structure that makes cross-plane hazard detection automatic in
our approach.  SVA monitors are typically simulation-only; synthesizable SVA
subsets exist but are not the norm.

Design-by-contract~\cite{meyer92} specifies preconditions, postconditions, and
invariants.  Our contribution is making the contracts algebraic theorems rather
than arbitrary predicates, encoding them as register bits, and proving that
fault states are terminal.

The oracle discipline is well-established in cryptography (independent test
vectors) and in instruction-set testing (Arm's ASL-to-RTL
cross-checking~\cite{reid16}).  To our knowledge, no prior hardware
verification methodology combines structurally different oracles with
typestate-based guards and poison-proof fault injection.

%=============================================================================
\section{Conclusion}
%=============================================================================

We have presented a three-layer hardware verification discipline combining
typestate analysis, structurally independent oracle testing, and poison-proof
fault injection.  The method is validated by a real bug ledger where each of
three independent layers caught a bug the others would have missed.  Five case
studies spanning arithmetic accelerators and a protocol state machine
demonstrate generality, and a measured head-to-head comparison with
SystemVerilog Assertions quantifies the tradeoffs.  All artifacts---testbenches,
synthesis scripts, oracle code, and this paper---are open-source under
permissive licenses.

%=============================================================================
\section*{Appendix --- Synthesis Metrics}
%=============================================================================

All metrics are post-route on Tang Primer 25K (Gowin GW5A-LV25MG121NES) using
Yosys + nextpnr-himbaechel + gowin\_pack.

\begin{table}[h]
\centering
\caption{Post-route resource utilization and Fmax}
\label{tab:synthesis}
\begin{tabular}{lcccccr}
\toprule
\textbf{Subsystem} & \textbf{LUT4} & \textbf{DFF} & \textbf{Fmax} & \textbf{DSP} &
\textbf{BSRAM} & \textbf{Si} \\
\midrule
$\varphi$ typestate guard$^1$ & 7 & 5 & --- & 0 & 0 & --- \\
$\varphi$ SVA guard$^1$ & 8 & 5 & --- & 0 & 0 & --- \\
IROTC probe & 1,670 & 532 & 59.9 & 0 & 0 & \checkmark \\
IROTC SPI (core) & 12,397 & 8,639 & 50.4$^2$ & 0 & 0 & bench \\
Lucas MAC probe & 697 & 216 & 119.0 & 0 & 0 & \checkmark \\
Lucas PHSLK probe & 293 & 146 & 200.4 & 0 & 0 & \checkmark \\
SPI link probe & 1,861 & 840 & 82.5 & 0 & 0 & \checkmark \\
RPLU2 arith probe & 9,211 & 7,926 & 54.4$^2$ & 0 & 0 & \checkmark \\
\bottomrule
\multicolumn{7}{l}{\footnotesize $^1$Post-synth only (not placed-and-routed).}\\
\multicolumn{7}{l}{\footnotesize $^2$Dual-clock-domain design (worst clock shown).}\\
\multicolumn{7}{l}{\footnotesize All rows are whole probe-top measurements from named}\\
\multicolumn{7}{l}{\footnotesize nextpnr logs (regenerable; see Appendix A of the}\\
\multicolumn{7}{l}{\footnotesize methods document for per-row evidence paths).}\\
\end{tabular}
\end{table}

%=============================================================================
\begin{thebibliography}{99}

\bibitem{strom86}
R.~E.~Strom and S.~Yemini,
``Typestate: A Programming Language Concept for Enhancing Software
Reliability,''
\textit{IEEE Trans.\ Softw.\ Eng.}, vol.~12, no.~1, pp.~157--171, 1986.

\bibitem{meyer92}
B.~Meyer,
``Applying `Design by Contract',''
\textit{Computer}, vol.~25, no.~10, pp.~40--51, 1992.

\bibitem{spu13}
J.~Curley,
``SPU-13 Sovereign Engine --- Architecture and Verification,''
GitHub repository, 2026.  \texttt{github.com/johncurley/SPU}

\bibitem{coconut24}
J.~Coconut,
``Typestate Analysis for Embedded Rust,''
\textit{Proc.\ EMSOFT}, 2024.

\bibitem{sva1800}
\textit{IEEE Standard for SystemVerilog}, IEEE Std 1800-2017.

\bibitem{reid16}
A.~Reid et al.,
``End-to-End Verification of ARM Processors with ISA-Formal,''
\textit{Proc.\ CAV}, 2016.

\bibitem{yanenko04}
E.~Yanenko,
``Evolving Categories: Consistent Framework for Representation of Data and
Algorithms,''
arXiv:cs/0406044, 2004.

\end{thebibliography}

\end{document}
